% !TeX root = main.tex
\section{Gerbes}
\label{sec:gerbes-and-twisted}
In this section, we collect some results from the theory on gerbes,
bandings by abelian groups and the relation to cohomology.
Throughout the section, $S$ will be an arbitrary algebraic stack.
By a \emph{stack} we mean a, not necessarily algebraic, stack in groupoids over the
big \emph{fppf} site of schemes over $S$. 
The symbol $\Delta$ will always denote an abelian group
in~$S_\fppf$. 
Later on, we will specialize to the situation when $\Delta$ is a diagonalizable group.
Accordingly, we will use multiplicative notation for
$\Delta$ and its cohomology groups $\HH^i(S_\fppf, \Delta)$.

\begin{definition}[{\cite[Definition~3.15]{lmb2000}, \cite[\sptag{06NZ}]{stacks-project}}]
A \emph{gerbe} over $S$ is a stack (in groupoids) $\alpha\colon X \to S$ satisfying the following
conditions.
\begin{enumerate}
\item
The diagonal $X \to X\times_S X$ is an epimorphism (see \cite[Definition~3.6]{lmb2000}).
\item
\label{def-gerbe-structure}
The structure morphism $\alpha\colon X \to S$ is an epimorphism.
\end{enumerate}
A gerbe is called \emph{trivial} if the structure morphism in \ref{def-gerbe-structure} splits.
\end{definition}

The prototypical example of a gerbe is the \emph{classifying stack} $\BB G$
for a group $G$ in $S_\fppf$.
In fact, any gerbe can be viewed as an \emph{fppf}-form of a classifying stack.

For any stack $\alpha\colon X \to S$,
we may consider its \emph{inertia stack} $I_{X/S} \to X$ (\cite[\sptag{034I}]{stacks-project}),
which is a group object in $X_\fppf$.
Its points over a morphism $x \colon U \to X$ are simply the 2-automorphisms $\gamma$ of $x$
mapping to the identity under the composition $\alpha \circ x$.
Any object $\mathcal{F}$ in $X_\fppf$ is endowed with a canonical right action by the inertia $I_{X/S}$,
called the \emph{inertial action}.
Explicitly, it is given on sections over $x$ by
\begin{equation}
\label{eq-inertial-action}
\mathcal{F}(x)\times I_{X/S}(x) \to \mathcal{F}(x),\qquad
(s, \gamma) \mapsto \mathcal{F}(\gamma)(s).
\end{equation}
We recall the following fundamental fact about the inertial action.
\begin{proposition}
\label{p-gerbe-pb-ff}
Let $\alpha\colon X \to S$ be a gerbe.
Then the functor $\alpha^{*}\colon S_\fppf \to X_\fppf$ is fully faithful with essential
image consisting of sheaves on which the inertia $I_{X/S}$ acts trivially.
\end{proposition}
\begin{proof}
See for instance \cite[Lemma~2.1.1.17]{lieblich2004}.
\end{proof}

\begin{definition}
A gerbe $\alpha\colon X \to S$ is called \emph{abelian} if the inertia group $I_{X/S} \to X$
is abelian.
We denote the full subcategory of stacks over $S$ which are abelian gerbes by
$\AbGerbe(S)$.
\end{definition}

The inertia acts on itself by conjugation.
In particular, Proposition~\ref{p-gerbe-pb-ff} shows that the inertia group of an abelian gerbe $\alpha\colon X \to S$
descends to an abelian
group object in $S_\fppf$, namely $\alpha_* I_{X/S}$.
Furthermore, given any 1-morphism $\rho\colon X \to Y$ of stacks over $S$,
we get an induced group homomorphism $I_{X/S} \to \rho^*I_{Y/S}$ in $X_\fppf$,
and if $X$ and $Y$ are abelian gerbes, this descends to a group homomorphism in $S_\fppf$.
This association gives a functor
\begin{equation}
\label{eq-band-functor}
\Band\colon \AbGerbe(S) \to \AbGroup(S)
\end{equation}
to the category of abelian group objects in $S_\fppf$.

\begin{definition}
\label{def-banding}
Let $\Delta$ be an abelian group object in $S_\fppf$.
The essential fiber of the functor \eqref{eq-band-functor} over $\Delta$ is called
the \emph{category of gerbes banded by $\Delta$} or \emph{of $\Delta$-gerbes}. 
Explicitly, a $\Delta$-gerbe is given by a pair $(\alpha, \iota)$ where $\alpha\colon X \to S$ is an abelian
gerbe and $\iota\colon \Delta \xrightarrow{\sim} \alpha_*I_{X/S}$ is a group isomorphism
called a \emph{banding} of $\alpha$ by $\Delta$.
The group $\Delta$ is called the band of $(\alpha,\iota)$.
A morphism of $\Delta$-gerbes is simply a morphism of stacks such that corresponding morphism on inertia
induces the identity on $\Delta$ via the bandings.
By abuse of notation, we usually denote a $\Delta$-gerbe $(\alpha, \iota)$ by $\alpha$ or even $X$.
\end{definition}

\begin{remark}
Giraud introduces a more general notion of banding applying to non-abelian gerbes
\cite[Chapter~IV]{giraud1971}.
We will not need this, more complicated, theory here.
\end{remark}

Note that a 1-morphism of gerbes over $S$ is always an epimorphism (\cite[Lemma~3.17]{lmb2000}),
and that it is an isomorphism if and only if it induces an isomorphism on inertia.
In particular, the category of $\Delta$-gerbes is a 2-groupoid.

\begin{example}
\label{ex:simple-gerbes}
Let $S$ be an algebraic stack,
and let $\Delta$ and $\Delta'$ be abelian groups in $S_\fppf$.
\begin{enumerate}
\item
The classifying stack $\BB \Delta$ is endowed with a canonical structure of $\Delta$-gerbe.
\item
Given a pair of abelian gerbes $X$ and $X'$ over $S$,
banded by $\Delta$ and $\Delta'$, respectively,
the product $X\times_S X'$ has a naturally defined banding by $\Delta\times \Delta'$.
\end{enumerate}
\end{example}

The following example illustrates that it is very natural to consider gerbes over
genuine algebraic stacks.
\begin{example}
\label{ex-gerbes}
Let $S$ be an algebraic stack. Assume that we have a central extension
$1 \to \Delta \to G \to H \to 1$ of groups in $S_\fppf$.
Then the induced morphism of classifying stacks $\BB G \to \BB H$
has a canonical structure of a $\Delta$-gerbe.
The following examples are two important special cases with $S = \Spec \mathbb{Z}$.
\begin{enumerate}
\item
\label{it-universal-projective}
The morphism $G \to H$ is the quotient map $\GL_n \to \PGL_n$.
In this case we obtain a $\GGm$-gerbe $\BB\GL_n \to
\BB\PGL_n$.
\item
\label{it-universal-root}
The morphism $G \to H$ is the morphism $\GGm \to \GGm$ given by
$x \mapsto x^n$. 
In this case the corresponding $\mu_n$-gerbe $\BB \GGm \to \BB \GGm$ can be thought
of as the $n$-th root stack of the universal line bundle $[\bA_1/\GGm]$ on $\BB
\GGm$ (cf.~\cite[Definition~2.2.6]{cadman2007}).
\end{enumerate}
\end{example}

Next we discuss functoriality of banded gerbes with respect to change of group.

\begin{construction}
\label{cons-gerbe-pf}
Let $\varphi\colon \Delta \to \Delta'$ be a homomorphism of abelian group objects in $S$,
and let $\alpha\colon X \to S$ be a gerbe banded by $\Delta$.
We construct a gerbe $\varphi_*\alpha\colon \varphi_*  X \to S$
banded by $\Delta'$ together with a morphism
\begin{equation}
\label{eq-gerbe-pf}
\rho \colon X \to \varphi_*X
\end{equation}
of gerbes over $S$ inducing the homomorphism $\varphi$ on the bands as follows.
 
Consider $X \times \BB \Delta'$ with its natural banding by $\Delta \times \Delta'$.
We define $\varphi_*X$ as the \emph{rigidification} (see \cite[Appendix~A]{aov2008})
of $X \times \BB \Delta'$ in the kernel of the epimorphism
$\Delta \times \Delta' \to \Delta'$ given by $(a, b) \mapsto \varphi(a) b$.
The morphism $\rho$ is the composition of the obvious map $X \to X\times \BB \Delta'$ followed by the
rigidification map.
\end{construction}

The construction above shows that given a gerbe banded by $\Delta$,
any morphism $\varphi\colon \Delta \to \Delta'$ lifts to a morphism of banded gerbes.
The following proposition,
which we surprisingly could not find in the literature,
shows that this lift is essentially unique.
In particular, Construction~\ref{cons-gerbe-pf} is functorial in a weak sense.

\begin{proposition}
\label{gerbe-functorial-group}
Let $\Delta$ be an abelian group in $S_\fppf$,
and let $\alpha\colon X \to S$ be a gerbe banded by $\Delta$.
Then the obvious functor
\begin{equation}
\Band_{X/} \colon \AbGerbe(S)_{X/} \to \AbGroup(S)_{\Delta/}
\end{equation}
from the category of abelian gerbes over $S$ under $X$ to the category of abelian groups in $S_\fppf$
under $\Delta$ is an equivalence of categories.
\end{proposition}
\begin{proof}
Construction~\ref{cons-gerbe-pf} shows that the functor
$\Band_{X/}$ is essentially surjective on objects.
Let $\rho\colon X \ra Y$ and  $\rho'\colon X \ra Y'$ be objects of $\AbGerbe(S)_{X/}$,
and denote their images in $\AbGroup(S)_{\Delta/}$ by
$\varphi\colon \Delta \to \Gamma$ and~$\varphi'\colon \Delta \to \Gamma'$, respectively.
We need to show that
\begin{equation}
\label{eq-band-equivalence}
\Hom_{X/}(\rho, \rho') \to \Hom_{\Delta/}(\varphi, \varphi')
\end{equation}
is and equivalence of categories.
To do so, we first note that \eqref{eq-band-equivalence} is functorial in~$S$,
and that we in fact have a morphism of stacks
\begin{equation}
\label{eq-band-stacky-equivalence}
\Phi\colon \sheafHom_{X/}(\rho, \rho') \to \sheafHom_{\Delta/}(\varphi, \varphi')
\end{equation}
over~$S$.
Explicitly, this morphism is constructed as the canonical map in the diagram
\begin{equation}
\label{eq-big-and-scary}
\xymatrix{
\sheafHom_{\Delta/}(\varphi, \varphi')\ar[ddd]\ar[rrr] & & & S \ar[ddd]^{\varphi'}\\
& \sheafHom_{X/}(\rho, \rho')\ar[d]\ar[r]\ar[ul]_\Phi & S\ar[ur]^\id \ar[d]^{\rho'} & \\
& \sheafHom(Y, Y') \ar[r]^{-\circ\rho}\ar[dl] & \sheafHom(X, Y')\ar[dr] & \\
\sheafHom(\Gamma, \Gamma') \ar[rrr]^{-\circ\varphi} & & & \sheafHom(\Delta, \Gamma'), \\
}
\end{equation}
where the inner and the outer squares are cartesian.

To verify that $\Phi$ is an isomorphism, we may work locally on~$S$.
In particular, we may assume that the stacks $X, Y$ and~$Y'$ are the classifying
stacks
$\BB \Delta$, $\BB \Gamma$ and~$\BB \Gamma'$, respectively.
Furthermore, we may assume that the morphisms $\rho$ and $\rho'$ are given by extension of torsors along
$\varphi$ and $\varphi'$, respectively.
In this situation, we have $\sheafHom(X, Y') \cong [\sheafHom(\Delta, \Gamma')/\Gamma']$,
where $\Gamma'$ acts trivially,
and the map to $\sheafHom(\Delta, \Gamma')$ is induced by forgetting the group action
(see e.g. Abramovich~\emph{et~al.}~\cite[Lemma~3.9(1)]{aov2008} and its proof).
Plugging this into \eqref{eq-big-and-scary},
together with the similar description for $\sheafHom(Y, Y')$ and the obvious descriptions for $\rho'$
and $-\circ\rho$,
it becomes clear that $\Phi$ is an isomorphism.
\end{proof}

\begin{construction}
\label{cons-gerbe-operations}
Let $\Delta$ be an abelian group in $S_\fppf$,
and let $\alpha, \beta\colon X \to S$ be $\Delta$-gerbes.
Denote the multiplication map on $\Delta$ by $m\colon \Delta \times \Delta \to \Delta$,
and the $n$-th power map, for $n\in \mathbb{Z}$, by~$p_n\colon \Delta \to \Delta$.
Using Construction~\ref{cons-gerbe-pf}, we define
\begin{equation}
\label{eq-gerbe-mult-inverse}
\alpha\beta := m_*(\alpha \times \beta), \qquad
\alpha^n := (p_n)_*\alpha.
\end{equation}
\end{construction}
Morally, Construction~\ref{cons-gerbe-operations} gives the 2-groupoid of $\Delta$-gerbes
the structure of an abelian group,
a statement which presumably could be made precise by use of Proposition~\ref{gerbe-functorial-group}.
At least, it is clear that the set of isomorphism classes of $\Delta$-gerbes forms an abelian group.

It is easy to see that we also have functoriality with respect to change of stacks.
That is, let $f\colon T \to S$ be a morphism of algebraic stacks
and let $\alpha\colon X \to S$ be a $\Delta$-gerbe.
Then the pull-back $f^*\alpha$ has a natural structure of $f^*\Delta$-gerbe.

\begin{theorem}[Giraud]
\label{t:gerbes-H2}
Let $S$ be an algebraic stack and let $\Delta$ be an abelian group in~$S_\fppf$.
Then the group $\HH^2(S_\fppf, \Delta)$ is canonically isomorphic to the set of isomorphism
classes of $\Delta$-gerbes over $S_\fppf$,
with its group structure induced by the operations in Construction~\ref{cons-gerbe-operations}.
Given a $\Delta$-gerbe $\alpha$ over $S_\fppf$,
we denote its class in $\HH^2(S_\fppf, \Delta)$ by~$[\alpha]$.
We have the following functorial properties:
\begin{enumerate}
\item
Given a morphism $f\colon T \to S$ of algebraic stacks,
we have $[f^*\alpha] = f^*[\alpha]$ in $H^2(T_\fppf, f^*\Delta)$.
\item
Given a homomorphism of abelian groups $\varphi\colon \Delta \to \Delta'$,
we have $[\varphi_\ast \alpha] = \varphi_*[\alpha]$ in $H^2(S_\fppf, \Delta')$.
\end{enumerate}
\end{theorem}
\begin{proof}
This is worked out in various places in \cite{giraud1971}.
Functoriality with respect to change of groups is described in Section~IV.3.1.
Comparison with usual cohomology is in Section~IV.3.4,
and functoriality with respect to change of topos in Section~V.1.
\end{proof}

%%% Local Variables: 
%%% mode: latex
%%% TeX-master: "main"
%%% End: 